% Clasificación exacta del entrelazador horizonte--TFD--Perron.

\section{Classification spectrale de l'entrelaceur entre horizon, purification thermochamp double et dynamique de Perron}
\label{sec:ivv-horizon-tfd-perron-classification}

Soient \(\mathcal H_{\mathrm H}\), \(\mathcal H_{\mathrm{TFD}}\) et
\(\mathcal H_{\mathrm P}\) les sous-espaces cycliques engendrés,
respectivement, par les microétats d'horizon, la purification
thermochamp double et l'opérateur positif de Perron. Notons
\(A_{\mathrm H}\), \(A_{\mathrm{TFD}}\) et \(A_{\mathrm P}\) leurs opérateurs
autoadjoints restreints et \(m_X(\lambda)\) la multiplicité spectrale
de \(\lambda\) dans le secteur \(X\).

\begin{theorem}[Existence et unicité à jauges près de l'entrelaceur spectral]
\label{thm:ivv-horizon-tfd-perron-classification}
Il existe une isométrie surjective
\[
 U:\mathcal H_{\mathrm H}\longrightarrow\mathcal H_{\mathrm{TFD}}
 \longrightarrow\mathcal H_{\mathrm P}
\]
qui entrelace les \(3\) opérateurs si et seulement si, après application des
normalisations déclarées d'énergie et de mémoire,
\[
 \operatorname{spec}(A_{\mathrm H})
 =\operatorname{spec}(A_{\mathrm{TFD}})
 =\operatorname{spec}(A_{\mathrm P}),
 \qquad
 m_{\mathrm H}(\lambda)=m_{\mathrm{TFD}}(\lambda)=m_{\mathrm P}(\lambda)
\]
pour tout \(\lambda\) du spectre commun. Lorsque ces conditions sont
satisfaites, l'entrelaceur est déterminé à une transformation unitaire près
dans chaque sous-espace propre de multiplicité supérieure à \(1\).
\end{theorem}

\begin{proof}
Si \(U A_X=A_YU\), l'équivalence unitaire préserve le spectre et les
multiplicités. Réciproquement, pour chaque \(\lambda\), on choisit des bases
orthonormées des sous-espaces propres de même dimension et l'on définit
\(U_\lambda\) par correspondance des bases. La somme orthogonale
\(U=\bigoplus_\lambda U_\lambda\) est unitaire et échange les opérateurs.
La liberté résiduelle est exactement le produit des groupes unitaires des
sous-espaces dégénérés.
\end{proof}

Le problème d'existence est ainsi clos par classification : la coïncidence
des spectres et des multiplicités produit l'entrelaceur ; son absence
démontre l'obstruction. Une égalité scalaire isolée ne remplace pas ces
conditions.

